\BourbakiExercisesSection

\begin{enumerate}
\def\labelenumi{\arabic{enumi}.}
\item
  在 § 8 习题 1 所确定的环结构中，哪些是域结构？
\item
  每个非零元素都可消去的有限环是域（参见 § 2，习题 6）。
\item
  设 G 为一个单的（§ 4，第 4 号，定义 7）带算子交换群。证明 G 的自同态环 A 是一个域。
\item
  设 K 为一个域。证明存在 K 的一个最小子域 F，并且 F 典范同构于域 Q 或域 Z/pZ（其中 p 为某个素数）。在前一种情形（相应地，后一种情形）下，称 K 的特征为 0（相应地，p）。
\item
  设 K 为一个特征 \(\neq2\) 的交换域；设 G 为 K 的加法群的一个子群，使得若 H 表示由 0 与 G 中元素 \(\neq0\) 的逆元组成的集合，则 H 也是 K 的加法群的一个子群。证明存在一个元素 \(a\in K\) 与 K 的一个子域 \(K'\) ，使得 \(G=aK'\) （首先证明：若 x 与 y 是 G 的元素且 \(y\neq0\) ，则 \(x^{2}/y\in G\) ；由此推出：若 x,y,z 是 G 的元素且 \(z\neq0\) ，则 xy/z\ensuremath{\in}G）。
\item
  设 K 为一个特征 \(\neq2\) 的交换域；设 f 为 K 到 K 中的一个映射，使得 \(f(x+y)=f(x)+f(y)\) 对所有 x 与 y 成立，且对所有 \(x\neq0\) 有 \(f(x)f(1/x)=1\) 。证明 f 是 K 到 K 的一个子域上的同构（证明 \(f(x^{2})=(f(x))^{2}\)）。
\item
  设 A 是一个交换环。
\end{enumerate}

\begin{enumerate}
\def\labelenumi{(\alph{enumi})}
\item
  每个素理想都是不可约的（§ 8，习题 11）。
\item
  素理想的集合对于关系 \(\subset\) 与 \(\supset\) 是归纳的。
\item
  设 \(\mathfrak{a}\) 为 \(\mathbf{A}\) 的一个异于 \(\mathbf{A}\) 的理想；令 \(\mathfrak{b}\) 为由这样的 \(x \in \mathbf{A}\) 组成的集合：存在一个整数 \(n \geqslant 0\) 使得 \(x^n \in \mathfrak{a}\) （\(n\) 依赖于 \(x\)）。证明 \(\mathfrak{b}\) 是一个理想，且等于 \(\mathbf{A}\) 的包含 \(\mathfrak{a}\) 的素理想之交。
\end{enumerate}

¶ 8. 环 A 称为布尔环，如果它的每个元素都是幂等的（换句话说，如果 \(x^2 = x\) 对所有 \(x \in \mathbf{A}\) 成立）。

\begin{enumerate}
\def\labelenumi{(\alph{enumi})}
\item
  在集合 \(\mathfrak{P}(\mathbf{E})\) （集合 \(\mathbf{E}\) 的子集组成的集合）中，证明通过令 \(\mathbf{AB} = \mathbf{A} \cap \mathbf{B}\) 与 \(\mathbf{A} + \mathbf{B} = (\mathbf{A} \cap \mathbb{C}\mathbf{B}) \cup (\mathbf{B} \cap \mathbb{C}\mathbf{A})\) 定义了一个布尔环结构。这个环同构于环 \(\mathbf{K}^{\mathrm{E}}\) ——\(\mathbf{E}\) 到模 2 的整数环 \(\mathbf{K} = \mathbf{Z}/(2)\) 中的映射所成的环（对每个子集 \(\mathbf{X} \subset \mathbf{E}\) ，考虑其“特征函数” \(\phi_{\mathbf{X}}\) ，满足 \(\phi_{\mathbf{X}}(x) = 1\) （若 \(x \in \mathbf{X}\)），\(\phi_{\mathbf{X}}(x) = 0\) （若 \(x \notin \mathbf{X}\)））。
\item
  每个布尔环 \(\mathbf{A}\) 都是交换的，且满足 \(x + x = 0\) （对 \(x \in \mathbf{A}\)）（把 \(x + x\) 写成幂等元，然后把 \(x + y\) 写成幂等元）。
\item
  如果布尔环 \(\mathbf{A}\) 不含零因子，则它或者化为 0，或者同构于 \(\mathbf{Z} / (2)\) （若 \(x\) 与 \(y\) 是 \(\mathbf{A}\) 的任意两个元素，证明 \(xy(x + y) = 0\)）。由此推出，在布尔环中每个素理想都是极大理想。
\item
  在布尔环 A 中，每个理想 \(a \neq A\) 都是包含 a 的素理想之交（应用习题 7(c)）。由此推出每个不可约理想都是极大理想（换句话说，在布尔环中，不可约理想、素理想与极大理想这三个概念是一致的）。
\item
  证明每个布尔环都同构于乘积环 \(\mathbf{K}^{\mathbb{E}}\) 的一个子环，其中 \(\mathbf{K} = \mathbf{Z}/(2)\) （利用 \((d)\)）。
\item
  若 \(p_{i}\) ( \(1 \leqslant i \leqslant n\) ) 是布尔环 A 中 n 个互不相同的极大理想，且 \(a = \bigcap_{1 \leqslant i \leqslant n} p_{i}\) ，证明商环 A/a 同构于乘积环 \(K^{n}\) 。由此推出每个有限布尔环都具有形式 \(K^{n}\) 。
\item
  在布尔环 A 中，关系 \(xy = x\) 是一个序关系；若把它记作 \(x \leqslant y\) ，证明 A 在此关系下是一个分配格（集合论，III，§ 1，习题 16），有最小元素 \(\alpha\) ，且对每个有序元素对 \((x, y)\) （满足 \(x \leqslant y\) ），存在一个元素 \(d(x, y)\) ，使得 \(\inf(x, d(x, y)) = \alpha\) , \(\sup(x, d(x, y)) = y\) 。反之，若一个有序集 A 具有这三条性质，证明合成律 \(xy = \inf(x, y)\) 与 \(x + y = d(\sup(x, y), \inf(x, y))\) 在 A 上定义了一个布尔环结构。
\end{enumerate}

\begin{enumerate}
\def\labelenumi{\arabic{enumi}.}
\setcounter{enumi}{8}
\tightlist
\item
  设 \(\mathbf{A}\) 为一个环，使得 \(x^3 = x\) 对所有 \(x \in \mathbf{A}\) 成立。我们要证明 \(\mathbf{A}\) 是交换的。
\end{enumerate}

\begin{enumerate}
\def\labelenumi{(\alph{enumi})}
\item
  证明 \(6\mathbf{A} = \{0\}\) ，且 \(2\mathbf{A}\) 与 \(3\mathbf{A}\) 是双边理想，满足 \(2\mathbf{A} + 3\mathbf{A} = \mathbf{A}\) 且 \(2\mathbf{A} \cap 3\mathbf{A} = \{0\}\) 。由此推出，可以假设或者 \(2\mathbf{A} = \{0\}\) ，或者 \(3\mathbf{A} = \{0\}\) 。
\item
  若 2A = \{0\}，计算 \((1 + x)^{3}\) ，推出 \(x^{2} = x\) 对所有 \(x \in A\) 成立，并利用习题 8 得出结论。
\item
  若 \(3\mathbf{A} = \{0\}\) ，计算 \((x + y)^3\) 与 \((x - y)^3\) ，证明 \(x^2y + xyx + yx^2 = 0\) 并用 \(x\) 左乘；由此推出 \(xy - yx = 0\) 。
\item
  设 A 为一个环，满足 \(3A = \{0\}\) 且 \(x^{3} = x\) 对所有 \(x \in A\) 成立；定义一个集合 I 以及 A 到 \((\mathbf{Z}/3\mathbf{Z})^{\mathrm{I}}\) 中的单射同态（方法与习题 8 相同）。
\end{enumerate}

\begin{enumerate}
\def\labelenumi{\arabic{enumi}.}
\setcounter{enumi}{9}
\tightlist
\item
  设 \(\mathbf{A}\) 是一个交换环， \(\mathfrak{U}\) 是其极大理想的交集。
\end{enumerate}

\begin{enumerate}
\def\labelenumi{(\alph{enumi})}
\item
  证明典范同态 \(\mathbf{A}^* \to (\mathbf{A} / \mathfrak{U})^*\) 是满射的，且其核为 \(1 + \mathfrak{U}\) 。
\item
  假设 A 的极大理想之集 M 是有限的，且 A* 是有限的。由此推出 A 是有限的（把 § 8 第 10 号的命题 8 应用于 A/U）。
\item
  由此推出素数集合是无限的（注意 \(\mathbf{Z}^{*} = \{1, -1\}\)）。
\item
  假设 A* 是有限的，且对每个极大理想 m，A/m 都是有限的。能否由此推出 A 是有限的？A 是可数的？
\end{enumerate}

\begin{enumerate}
\def\labelenumi{\arabic{enumi}.}
\setcounter{enumi}{10}
\tightlist
\item
  设 P 为素数集合，A 为域 \(\mathbf{Z} / p\mathbf{Z}\) , \(p \in \mathbb{P}\) 的乘积环。设 \(a\) 为 A 的由元素 \((a_p)_{p \in \mathbb{P}}\) （满足 \(a_p \neq 0\) 仅对有限多个指标 \(p\) 成立）组成的子集。
\end{enumerate}

\begin{enumerate}
\def\labelenumi{(\alph{enumi})}
\item
  a 是 A 的一个理想。
\item
  设 \(\mathbf{B} = \mathbf{A} / a\) 。对每个整数 \(n > 0\) 与每个元素 \(b \neq 0\) （在 \(\mathbf{B}\) 中），存在且仅存在一个元素 \(b'\) （属于 \(\mathbf{B}\)），使得 \(nb' = b\) （注意，若 \(p \in \mathbb{P}\) 不整除 \(n\) ，则用 \(n\) 乘在 \(\mathbf{Z} / p\mathbf{Z}\) 中是双射）。由此推出 \(\mathbf{B}\) 包含一个同构于 \(\mathbf{Q}\) 的子域。
\end{enumerate}

\begin{enumerate}
\def\labelenumi{\arabic{enumi}.}
\setcounter{enumi}{11}
\item
  设 A 为一个交换环，且 \(a, b \in A\) 。证明 ab 在 \(\mathrm{A}/(a - a^{2}b)\) 中的典范像是幂等的。给出一个使该幂等元异于 0 与 1 的例子。
\item
  确定环 Z 与环 \(Z \times Z\) 的自同态。更一般地，若 I 与 J 是有限集，确定环 \(Z^{1}\) 到环 \(Z^{J}\) 中的同态。
\item
  设 A 与 B 是两个环，f 与 g 是 A 到 B 中的两个同态。使 \(x \in A\) 满足 \(f(x) = g(x)\) 者组成的集合是 A 的子环、A 的理想吗？
\end{enumerate}

¶ 15. 设 A 是一个无零因子的非交换环。称 A 容许一个左分式域，如果它同构于某个域 K 的子环 B，使得 K 的每个元素都具有形式 \(x^{-1}y\) ，其中 \(x \in \mathbf{B}, y \in \mathbf{B}\) 。

\begin{enumerate}
\def\labelenumi{(\alph{enumi})}
\tightlist
\item
  设 A' 为 A 中元素 \(\neq0\) 组成的集合。为使 A 容许一个左分式域，下列条件必须满足：
\end{enumerate}

\begin{enumerate}
\def\labelenumi{(\Alph{enumi})}
\setcounter{enumi}{6}
\tightlist
\item
  对所有 \(x \in \mathbf{A}\) , \(x' \in \mathbf{A}'\) ，存在 \(u \in \mathbf{A}'\) 与 \(v \in \mathbf{A}\) ，使得 \(ux = vx'\) 。
\end{enumerate}

\begin{enumerate}
\def\labelenumi{(\alph{enumi})}
\setcounter{enumi}{1}
\tightlist
\item
  反之，假设条件 (G) 成立。证明在集合 A × A' 中，\((x, x')\) 与 \((y, y')\) 之间的关系 R——即“对每个由非零元素组成的有序对 \((u, v)\) ，只要 \(ux' = vy'\) ，就有 \(ux = vy\) ”——是一个等价关系。
\end{enumerate}

设 \((x, x')\) , \((y, y')\) 为 \(A \times A'\) 的两个元素，\(\xi\) 与 \(\eta\) 为它们各自的类（mod. R）。对每个有序对 \((u, u') \in A \times A'\) （满足 \(u'x = uy'\) ），证明 \((uy, u'y')\) 的类（mod. R）仅依赖于类 \(\xi\) 与 \(\eta\) ；若把它记作 \(\xi\eta\) ，就在集合 \(K = (A \times A')/R\) 上定义了一个合成律。若 \(K'\) 为 K 中异于诸元素 \((0, x')\) （属于 \(A \times A'\)）的类 0 的元素组成的集合，则 \(K'\) 连同由上述法则诱导的法则构成一个群。

对每个元素 \(x \in \mathbf{A}\) ，诸元素 \((x'x, x')\) （其中 \(x'\) 遍历 \(\mathbf{A}'\) ）定义了 \(\mathbf{A}\) （仅带乘法）到 \(\mathbf{K}\) 的一个子环上的同构。把 \(\mathbf{A}\) 与它在此同构下的像等同以后，有序对 \((x, x') \in \mathbf{A} \times \mathbf{A}'\) 的类（mod. R）就与元素 \(x'^{-1}x\) 等同。

既然如此，若 \(\xi = x'^{-1}x\) 是 K 的一个元素，且 1 表示 \(\mathbf{K}'\) 的单位元素，则 \(\xi + 1\) 表示元素 \(x'^{-1}(x + x')\) ，它不依赖于 \(\xi\) 写成形式 \(x'^{-1}x\) 的方式。然后我们写 \(\xi + 0 = \xi\) ，且对 \(\eta \neq 0\) ，写 \(\xi + \eta = \eta(\eta^{-1}\xi + 1)\) 。证明这样在 K 上定义的加法与乘法在该集合上确定了一个域结构，它扩充了 A 上的环结构；换言之，条件 (G) 使 A 容许一个左分式域是充分的。

\begin{enumerate}
\def\labelenumi{\arabic{enumi}.}
\setcounter{enumi}{15}
\item
  设 A 是一个环，其中每个非零元素都是可消去的。为使 A 容许一个左分式域（习题 15），必要且充分的是：在 A 中，两个异于 0 的左理想之交永不缩减为 0。
\item
  设 \((\mathbf{K}_i)_{i\in I}\) 是一个域族。对每个子集 \(J\) ⊂ \(I\) ，设 \(e_J\) 表示乘积 \(\mathbf{K} = \prod_{i\in I}\mathbf{K}_i\) 中这样的元素：它的第 \(i\) 个分量当 \(i\in J\) 时等于 0，而当 \(i\in I - J\) 时等于 1。
\end{enumerate}

\begin{enumerate}
\def\labelenumi{(\alph{enumi})}
\item
  设 \(x = (x_{i})\) 是 \(\mathbf{K}\) 的一个元素，\(J(x)\) 为元素 \(i \in I\) 中使 \(x_{i} = 0\) 者组成的集合。证明存在一个可逆元素 \(u\) （在 \(\mathbf{K}\) 中），使得 \(x = ue_{J(x)} = e_{J(x)}u\) 。
\item
  设 \(\mathfrak{a}\) 是 \(\mathbf{K}\) 的一个异于 \(\mathbf{K}\) 的左（相应地，右）理想。设 \(\mathfrak{F}_{\mathfrak{a}}\) 为子集 \(\mathbf{J}\) （属于 \(\mathbf{I}\) 且满足 \(e_{J} \in \mathfrak{a}\)）组成的集合。证明 \(\mathfrak{F}_{\mathfrak{a}}\) 具有下列性质：
\end{enumerate}

\((b_{1})\) 若 \(\mathbf{J} \in \mathfrak{F}_{\mathfrak{a}}\) 且 \(\mathbf{J}' \supset \mathbf{J}\) ，则 \(\mathbf{J}' \in \mathfrak{F}_{\mathfrak{a}}\) 。

\((b_{2})\) 若 \(\mathrm{J}_1, \mathrm{J}_2 \in \mathfrak{F}_{\mathfrak{a}}\) ，则 \(\mathrm{J}_1 \cap \mathrm{J}_2 \in \mathfrak{F}_{\mathfrak{a}}\) 。

\[
\left(b _ {3}\right) \varnothing \notin \mathfrak {F} _ {\mathfrak {a}}.
\]

证明 \(x \in \mathfrak{a}\) 等价于 \(J(x) \in \mathfrak{F}_{\mathfrak{a}}\) （利用 (a)）。由此推出 \(\mathfrak{a}\) 是一个双边理想。

\begin{enumerate}
\def\labelenumi{(\alph{enumi})}
\setcounter{enumi}{2}
\item
  I 的一个子集族称为滤子，如果它满足上述性质 \((b_{1})\) , \((b_{2})\) , \((b_{3})\) （参见《一般拓扑学》I，§6）。证明映射 \(\mathfrak{a} \mapsto \mathfrak{F}_{\mathfrak{a}}\) 是 K 的异于 K 的理想之集到 I 的滤子之集上的一个严格递增双射。*证明 a 是极大的当且仅当 \(\mathfrak{F}_{\mathfrak{a}}\) 是一个超滤子。*
\item
  *设 \(\mathfrak{F}\) 是素数集合 \(P\) 上的一个非平凡超滤子，\(a\) 为环 \(\mathbf{K} = \prod_{p \in P} \mathbf{Z}/p\mathbf{Z}\) 的相应理想，并设 \(k = K/a\) 。证明 \(k\) 是一个特征为 \(0\) 的域。*
\end{enumerate}

¶ 18. (a) 设 K 是一个域，a、b 是 K 中两个不可交换的元素。证明

\[
\begin{array}{l} (1) a = (b - (a - 1) ^ {- 1} b (a - 1)) (a ^ {- 1} b a - (a - 1) ^ {- 1} b (a - 1)) ^ {- 1} \\ (2) a = (1 - (a - 1) ^ {- 1} b ^ {- 1} (a - 1) b) (a ^ {- 1} b ^ {- 1} a b - (a - 1) ^ {- 1} b ^ {- 1} (a - 1) b) ^ {- 1}. \end{array}
\]

\begin{enumerate}
\def\labelenumi{(\alph{enumi})}
\setcounter{enumi}{1}
\item
  设 K 是一个非交换域，x 是不属于 K 的中心的一个元素。证明 K 由 x 的共轭元素 \(axa^{-1}\) 组成的集合生成。（设 \(K_{1}\) 为 K 的由 x 的共轭元素集合生成的子域。由 (a) 推出，若 \(K_{1} \neq K\) 且 \(a \in K_{1}\) 而 \(b \notin K_{1}\) ，则必有 ab = ba。通过在 \(K_{1}\) 中考虑两个不可交换的元素 a、\(a'\) 以及一个元素 \(b \notin K_{1}\) ，并注意到 \(ba \notin K_{1}\) ，由此得出矛盾。）
\item
  由 (b) 推出，若 Z 是 K 的中心且 H 是 K 的一个子域，使得 \(Z \subset H \subset K\) 且 \(aHa^{-1} = H\) 对 K 中所有 \(a \neq 0\) 成立，则 H = Z 或 H = K（嘉当–布劳尔–华定理）。
\item
  由 (a) 推出，在一个非交换域 K 中，换位子 \(aba^{-1}b^{-1}\) （元素 \(\neq0\) 的换位子）所组成的集合生成域 K。
\end{enumerate}

\begin{enumerate}
\def\labelenumi{\arabic{enumi}.}
\setcounter{enumi}{18}
\tightlist
\item
  设 A 是一个非零伪环，使得对 A 中所有 \(a \neq 0\) 及所有 \(b \in \mathbf{A}\) ，方程 \(ax + ya = b\) 都有由 A 的元素组成的有序对 \((x, y)\) 作为解；换言之，\(a\mathbf{A} + \mathbf{A}a = \mathbf{A}\) 对 A 中的 \(a \neq 0\) 成立。
\end{enumerate}

\begin{enumerate}
\def\labelenumi{(\alph{enumi})}
\tightlist
\item
  证明 A 是 A 中唯一的非零双边理想。
\end{enumerate}

\begin{enumerate}
\def\labelenumi{(\alph{enumi})}
\setcounter{enumi}{1}
\item
  证明关系 \(a \neq 0\) 蕴含 \(a^2 \neq 0\) （注意，若 \(a^2 = 0\) ，则将得出 \(aAa = \{0\}\) ）。
\item
  证明：若 \(ab = 0\) 且 \(a \neq 0\) （相应地，\(b \neq 0\) ），则 \(b = 0\) （相应地，\(a = 0\) ）。（注意 \((ba)^2 = 0\) ，并由此推出 \(bAb = \{0\}\) 。）
\end{enumerate}

\begin{enumerate}
\def\labelenumi{\arabic{enumi}.}
\setcounter{enumi}{19}
\item
  设 A 是一个非零伪环，使得存在 \(a \in A\) ，对每个 \(b \in A\) ，方程 ax = b、xa = b 中有一个具有解 \(x \in A\) 。证明：若 A 除 0 外没有零因子，则 A 具有单位元素。
\item
  设 A 是一个非零伪环，其中对所有 \(a \neq 0\) 以及所有 \(b \in A\) ，方程 ax = b、xa = b 中有一个具有解 \(x \in A\) 。证明 A 是一个域（利用习题 19 和 20）。
\end{enumerate}

